\chapter{Mediterranean diet}

\section*{Definition and Overview}
The Mediterranean diet is a nutritional framework inspired by the traditional dietary patterns of populations living in the Mediterranean basin, particularly Greece, southern Italy, and Spain, during the mid-20th century. It is characterised by a high consumption of plant-based foods, including vegetables, fruits, whole grains, legumes, nuts, and seeds; a moderate-to-high intake of monounsaturated fats (principally extra virgin olive oil); a moderate consumption of fish and poultry; a low-to-moderate intake of dairy products (mostly cheese and yoghurt); and a low intake of red meat, processed meats, and refined sugars. Moderate consumption of red wine, typically with meals, is also a traditional component. In modern medicine, the Mediterranean diet is recognised as one of the healthiest dietary patterns globally, backed by robust clinical evidence demonstrating its efficacy in reducing cardiovascular disease, type 2 diabetes, cognitive decline, and overall mortality.

\section*{Epidemiology}
The health benefits of the Mediterranean diet were first documented in the landmark Seven Countries Study in the 1950s, which observed significantly lower rates of coronary heart disease and all-cause mortality in Mediterranean cohorts compared to Northern European and North American cohorts, despite a high total fat intake. In Australia, chronic non-communicable diseases, such as cardiovascular disease and type 2 diabetes, represent the leading causes of death and disability. Large-scale clinical trials and prospective cohort studies have demonstrated that adherence to a Mediterranean-style diet reduces the risk of major cardiovascular events by approximately 30\% in high-risk individuals. The diet is recommended by major international and Australian health organisations, including the National Heart Foundation of Australia, as a primary and secondary prevention strategy for metabolic and vascular diseases.

\section*{Aetiology and Pathophysiology}
The clinical efficacy of the Mediterranean diet lies in its synergy of nutrient-dense components that target the pathophysiological mechanisms of chronic cardiovascular and metabolic diseases.
\begin{itemize}
    \item \textbf{Anti-inflammatory and antioxidant effects:} The diet is exceptionally rich in polyphenols, carotenoids, vitamins C and E, and selenium, which scavenge free radicals and reduce systemic oxidative stress and chronic inflammation (measured by lower C-reactive protein levels).
    \item \textbf{Lipid profile optimisation:} Monounsaturated fatty acids (primarily oleic acid from extra virgin olive oil) and polyunsaturated fatty acids (from nuts and seeds) lower low-density lipoprotein cholesterol and triglycerides while maintaining or increasing high-density lipoprotein cholesterol.
    \item \textbf{Endothelial function improvement:} High intake of omega-3 fatty acids (specifically eicosapentaenoic acid and docosahexaenoic acid from fatty fish) promotes nitric oxide synthesis in endothelial cells, leading to vasodilation, improved arterial compliance, and blood pressure reduction.
    \item \textbf{Insulin sensitivity enhancement:} The high dietary fibre content from legumes, vegetables, and whole grains slows carbohydrate digestion and absorption, reducing postprandial glucose excursions and improving pancreatic beta-cell function.
    \item \textbf{Gut microbiome modulation:} Soluble and insoluble fibres act as prebiotics, promoting the growth of beneficial gut bacteria that produce short-chain fatty acids, which have immunomodulatory and metabolic benefits.
\end{itemize}

\section*{Clinical Features}
The clinical presentation that warrants the adoption of a Mediterranean diet includes metabolic risk factors, established cardiovascular disease, or a desire for general health promotion.
\begin{itemize}
    \item Hypercholesterolaemia or dyslipidaemia, characterised by elevated total cholesterol and LDL cholesterol, and low HDL cholesterol.
    \item Hypertension, defined as persistent blood pressure readings of 140/90 mmHg or higher.
    \item Impaired fasting glucose, insulin resistance, or established Type 2 diabetes mellitus.
    \item Metabolic syndrome, presenting as a combination of abdominal obesity, hypertension, dyslipidaemia, and impaired glucose tolerance.
    \item Established cardiovascular disease, including coronary artery disease, peripheral arterial disease, or a history of stroke.
    \item Overweight or obesity, particularly abdominal (visceral) adiposity.
    \item Mild cognitive impairment or early-stage neurodegenerative diseases, where diet plays a role in slowing progression.
\end{itemize}

\section*{Differential Diagnosis}
In a dietary and metabolic context, the Mediterranean diet must be compared and distinguished from other popular nutritional patterns:
\begin{itemize}
    \item \textbf{Ketogenic diet:} A high-fat, very low-carbohydrate diet that induces ketosis. While it can cause rapid weight loss, it restricts whole grains, fruits, and legumes, and may elevate LDL cholesterol, unlike the cardioprotective Mediterranean diet.
    \item \textbf{Low-fat diets:} Traditionally recommended for cardiovascular health, but clinical trials have shown that the Mediterranean diet (rich in healthy fats) is superior in reducing cardiovascular events and improving lipid profiles.
    \item \textbf{Paleo diet:} Focuses on unprocessed meats, fish, vegetables, and nuts, but excludes grains, legumes, and dairy. The Mediterranean diet, by contrast, includes whole grains and legumes, which are rich in cardioprotective soluble fibres.
    \item \textbf{Western diet:} The standard modern diet high in refined grains, added sugars, trans fats, and red/processed meats; the primary dietary driver of global metabolic diseases, directly contrasted with the Mediterranean diet.
\end{itemize}

\section*{When to Seek Medical Care}
Individuals should consult their doctor or an Accredited Practising Dietitian before making major changes to their diet, especially if they have chronic health conditions.

\textbf{Call 000 (or local emergency services) for:}
\begin{itemize}
    \item Sudden, severe chest pain, tightness, or pressure radiating to the arm, neck, or jaw (signs of an acute myocardial infarction).
    \item Sudden weakness or numbness in the face, arm, or leg, difficulty speaking, or loss of balance (signs of an acute stroke).
    \item Sudden, severe shortness of breath.
\end{itemize}

Other non-emergency reasons to consult a healthcare provider include:
\begin{itemize}
    \item Wishing to transition to a Mediterranean diet to manage high cholesterol, high blood pressure, or pre-diabetes.
    \item Requiring assistance in tailoring the diet to manage specific medical conditions (such as chronic kidney disease, which requires monitoring of potassium and phosphate).
    \item Experiencing unintended weight loss or digestive issues after changing eating habits.
    \item Seeking advice on how to safely manage diabetes medications (such as insulin) when increasing dietary fibre and reducing processed carbohydrates.
\end{itemize}

\section*{Diagnosis and Investigations}
Diagnostic evaluation prior to and during the adoption of the Mediterranean diet involves checking cardiovascular and metabolic biomarkers.
\begin{itemize}
    \item \textbf{Fasting lipid profile:} Measures total cholesterol, LDL-C, HDL-C, and triglycerides to establish baseline cardiovascular risk and monitor diet efficacy.
    \item \textbf{Glycated haemoglobin:} Assesses long-term blood glucose control (HbA1c) to screen for or monitor type 2 diabetes.
    \item \textbf{Inflammatory markers:} High-sensitivity C-reactive protein test to evaluate systemic cardiovascular inflammation.
    \item \textbf{Blood pressure monitoring:} Regular checks to assess response to dietary sodium reduction and increased potassium intake.
    \item \textbf{Body composition analysis:} Measuring body weight, body mass index, and waist circumference to track visceral fat reduction.
    \item \textbf{Nutritional assessment:} Review of dietary intake by an Accredited Practising Dietitian to identify nutrient gaps and guide food choices.
\end{itemize}

\section*{Management}
The implementation of the Mediterranean diet involves a shift in food choices, cooking habits, and lifestyle practices.
\begin{itemize}
    \item \textbf{Primary fat source:} Replacing butter, margarine, and refined oils with extra virgin olive oil for cooking, baking, and dressing salads.
    \item \textbf{Abundant plant consumption:} Eating at least 5--7 servings of vegetables and 2--3 servings of fresh fruit daily, incorporating them into every meal.
    \item \textbf{Whole grains and legumes:} Choosing whole-grain bread, oats, brown rice, quinoa, and pasta, and eating legumes (chickpeas, lentils, beans) at least 3--4 times a week.
    \item \textbf{Fatty fish intake:} Eating oily fish (such as salmon, sardines, mackerel, or tuna) at least twice a week to ensure adequate omega-3 fatty acid intake.
    \item \textbf{Nuts and seeds:} Consuming a small handful (approximately 30 grams) of unsalted nuts (walnuts, almonds) and seeds daily as snacks or salad toppings.
    \item \textbf{Limiting red and processed meats:} Restricting red meat (beef, lamb) to a few times a month, and avoiding processed meats (sausages, bacon, deli meats) completely.
    \item \textbf{Moderate dairy and poultry:} Consuming moderate portions of poultry and fermented dairy products, such as Greek yoghurt and cheese, weekly.
    \item \textbf{Lifestyle integration:} Emphasising physical activity, sharing meals with family and friends (social connection), and enjoying food mindfully.
\end{itemize}

\section*{Complications}
There are no medical complications associated with a balanced Mediterranean diet; however, potential challenges or misapplications can occur.
\begin{itemize}
    \item \textbf{Caloric excess from healthy fats:} Overconsuming olive oil, nuts, and seeds can lead to excessive calorie intake and weight gain if portions are not managed.
    \item \textbf{Inadequate calcium intake:} If dairy products are overly restricted without replacing them with fortified plant milks or calcium-rich vegetables.
    \item \textbf{Medication interactions:} High intake of leafy green vegetables rich in vitamin K (such as spinach or kale) can interfere with the anticoagulant effect of warfarin, requiring stable intake and INR monitoring.
    \item \textbf{Digestive discomfort:} Rapidly increasing dietary fibre without adequate fluid intake can cause transient bloating, gas, and abdominal cramps.
    \item \textbf{Risks of excessive alcohol:} Excessive red wine consumption can lead to liver disease, increased cancer risk, and dependency; intake must remain moderate (no more than one glass daily for women, two for men).
\end{itemize}

\section*{Prognosis and Outcomes}
Adopting a Mediterranean diet carries a highly favourable clinical prognosis. Long-term studies indicate that high adherence to the diet is associated with a 20--30\% reduction in cardiovascular mortality, a lower risk of developing type 2 diabetes, better glycaemic control in existing diabetes, and a slower rate of age-related cognitive decline. Furthermore, it is highly sustainable over the long term compared to restrictive diets, leading to permanent, positive lifestyle improvements and enhanced quality of life.

\section*{Prevention and Self-Care}
\begin{itemize}
    \item Cook at home more often using fresh, whole ingredients and extra virgin olive oil as your primary fat.
    \item Keep a bowl of fresh fruit on the counter as a primary snack, replacing processed biscuits, chips, or sweets.
    \item Add a serving of vegetables to every meal; for example, toss spinach into your morning eggs, or add a side salad to your lunch.
    \item Incorporate meat-free days into your weekly routine, using lentils, chickpeas, or beans as your protein source.
    \item Choose unsalted, raw, or dry-roasted nuts and seeds instead of processed snacks.
    \item Stay hydrated by choosing water as your primary beverage, and limit sugar-sweetened drinks and juices.
    \item If you choose to drink alcohol, limit red wine to a single glass with your main meal, and avoid binge drinking.
    \item Enjoy your meals slowly and mindfully, ideally sitting down with family or friends to foster social connection.
\end{itemize}

\section*{Sources and Further Reading}
The following organisations provide authoritative, evidence-based information on the Mediterranean diet and cardiovascular health:
\begin{itemize}
    \item National Heart Foundation of Australia. \textit{Healthy eating guide and Mediterranean diet recommendations.} https://www.heartfoundation.org.au/
    \item Dietitians Australia. \textit{Mediterranean diet advice from Accredited Practising Dietitians.} https://dietitiansaustralia.org.au/
    \item Mayo Clinic. \textit{Mediterranean diet: a heart-healthy eating plan.} https://www.mayoclinic.org/healthy-lifestyle/nutrition-and-healthy-eating/in-depth/mediterranean-diet/art-20047801
    \item Harvard T.H. Chan School of Public Health. \textit{Diet review: Mediterranean diet.} https://www.hsph.harvard.edu/nutritionsource/healthy-weight/diet-reviews/mediterranean-diet/
    \item NHS. \textit{What is a Mediterranean-style diet?} https://www.nhs.uk/live-well/eat-well/how-to-eat-a-balanced-diet/what-is-a-mediterranean-diet/
\end{itemize}
